% called by main.tex
%

\section*{Abstract}
\addcontentsline{toc}{section}{Abstract}
\label{sec::abstract}

Large language models generate text as flat token sequences, but many technical artefacts are only correct when they preserve an underlying hierarchy. This mismatch is especially visible in Kubernetes manifests: although they are serialized as YAML text, their validity depends on nested mappings, lists, indentation levels, required fields and relations between resources. This thesis studies that tension through the automatic generation of Kubernetes manifests from natural language instructions, with the goal of evaluating not only whether the generated output resembles YAML, but whether it can be reconstructed and validated as a structured object.

The work defines a parser-facing intermediate representation based on ordered line blocks with an explicit hierarchical level. From this representation, a deterministic parser reconstructs the final YAML manifest without adding hidden repair logic. On top of this contract, the thesis compares several modeling strategies. The baseline uses the base instruction model without supervised adaptation. The first supervised branch, \texttt{serialized\_sft}, asks the model to emit both line text and level as part of the generated surface. The second branch, \textit{Two-Head Model}, separates the problem by adding an explicit structural head for level prediction. Additional ordinal and positional variants explore whether the hierarchical signal can be improved, and a lightweight DPO stage studies whether automatic preference optimization can refine a supervised model once the structural format is stable.

The results show that supervised serialization of the hierarchy produces the strongest structural improvement over the baseline. In the final test comparison, \texttt{serialized\_sft} reaches 97.14 percent YAML parseability, compared with 46.77 percent for the baseline, and also improves line fidelity, level accuracy, prompt adequacy and approximate semantic coverage. The explicit two-head architecture improves over the baseline in several metrics, but in its current form it does not match the robustness of the serialized branch, mainly because text generation, level prediction and parser reconstruction must remain tightly coordinated. DPO is useful when applied after a stable supervised policy: the best DPO variant reaches 100 percent YAML parseability in test and obtains the highest Kubernetes domain score among the evaluated models.

The main conclusion is therefore nuanced. Making structure explicit is useful, but not every explicit structural mechanism improves the final system automatically. In this project, the most reliable path is to expose hierarchy in a supervised block representation, reconstruct YAML through a deterministic parser and evaluate the result with syntax, structure, prompt and domain-oriented metrics. The work also shows that surface similarity is not enough for this kind of task: the relevant question is whether a flat autoregressive generator can produce an output whose hierarchical organization remains coherent after reconstruction.


\newpage
\section*{Resumen}
\addcontentsline{toc}{section}{Resumen}
\label{sec::resumen}


Los modelos de lenguaje generan texto como una secuencia plana de tokens, pero muchos artefactos técnicos solo son correctos cuando conservan una estructura jerárquica. Esta tensión aparece con claridad en los manifiestos de Kubernetes: aunque se escriben como YAML, su validez depende de mapas anidados, listas, niveles de indentación, campos obligatorios y relaciones entre recursos. Esta memoria estudia ese problema a partir de la generación automática de manifiestos Kubernetes desde instrucciones en lenguaje natural, con el objetivo de evaluar no solo si la salida se parece a YAML, sino si puede reconstruirse y validarse como un objeto estructurado.

El trabajo define una representación intermedia orientada al parser, basada en bloques de línea ordenados y acompañados de un nivel jerárquico explícito. A partir de esa representación, un parser determinista reconstruye el manifiesto YAML final sin introducir reparaciones ocultas. Sobre este contrato se comparan varias estrategias de modelado. La línea base utiliza el modelo de instrucciones sin ajuste supervisado. La primera rama supervisada, \texttt{serialized\_sft}, hace que el modelo genere tanto el texto de cada línea como su nivel en la propia superficie de salida. La segunda rama, \textit{Two-Head Model}, separa ambas señales mediante un cabezal estructural explícito para predecir \texttt{level}. Además, se exploran variantes ordinales y posicionales para estudiar la predicción jerárquica, y una etapa ligera de DPO para comprobar si la optimización por preferencias automáticas puede refinar un modelo supervisado ya estable.

Los resultados muestran que la serialización supervisada de la jerarquía produce la mejora estructural más clara respecto a la línea base. En la comparación final sobre test, \texttt{serialized\_sft} alcanza un 97,14\,\% de YAML parseable frente al 46,77\,\% del baseline, y mejora también la fidelidad de líneas, la exactitud de \texttt{level}, la adecuación al prompt y la cobertura semántica aproximada. La arquitectura con cabezal explícito mejora a la línea base en varias métricas, pero en su configuración actual no alcanza la robustez global de la rama serializada, principalmente porque el texto generado, el nivel predicho y la reconstrucción del parser deben mantenerse muy coordinados. DPO resulta útil cuando parte de una política supervisada estable: la mejor variante alcanza un 100\,\% de YAML parseable en test y obtiene el mejor score de dominio Kubernetes entre los modelos evaluados.

La conclusión principal es, por tanto, matizada. Hacer explícita la estructura ayuda, pero no cualquier mecanismo estructural explícito mejora automáticamente el sistema final. En este proyecto, la vía más fiable consiste en exponer la jerarquía mediante una representación supervisada por bloques, reconstruir el YAML con un parser determinista y evaluar la salida con métricas sintácticas, estructurales, de adecuación al prompt y de dominio. El trabajo muestra además que la similitud superficial no basta para este tipo de tarea: la pregunta relevante es si un generador autoregresivo plano puede producir una salida cuya organización jerárquica siga siendo coherente después de la reconstrucción.
